\chapter{परिमाण की कोटि: ब्रह्मांड को मापना}\label{ch:g10:orders-of-magnitude}

भौतिकी वहीं से आरंभ होती है जहाँ राय समाप्त होती है: माप से। पर भौतिकी
जिन चीज़ों को मापती है, उनका फैलाव बेतुका है --- एक परमाणु नाभिक लगभग
\qty{e-15}{m} चौड़ा है, प्रेक्षणीय ब्रह्मांड लगभग \qty{e27}{m}: $10^{42}$ का
गुणक, जिसे कोई मानसिक चित्र सँभाल ही नहीं सकता। यह अध्याय वही औज़ार
देता है जो इस फैलाव को वश में करते हैं --- मात्रक, दस की घातें,
\omterm{def:g10:orders-of-magnitude:sigfig}{सार्थक अंक} --- और दो आदतें, जो इसके बाद हर पृष्ठ पर काम आएँगी:
मात्रक लिखिए, और \omterm{def:g10:orders-of-magnitude:oom}{परिमाण की कोटि} जाँचिए।

\section{मात्रक: मापन का व्याकरण}

\begin{definition}[भौतिक राशियाँ और SI मात्रक]\label{def:g10:orders-of-magnitude:unit}
\emph{भौतिक राशि}\index{भौतिक राशि} संसार का वह गुण है जिसे मापा जा
सकता है: कोई लंबाई, कोई अवधि, कोई द्रव्यमान। उसे मापने का अर्थ है यह
गिनना कि कोई संदर्भ राशि --- \emph{मात्रक}\index{मात्रक} --- उसमें
कितनी बार समाती है। \emph{मात्रकों की अंतर्राष्ट्रीय
प्रणाली}\index{मात्रकों की अंतर्राष्ट्रीय प्रणाली (SI)} (SI) सात मूल
मात्रक तय करती है, जिनमें से तीन इस पूरे खंड को चलाते हैं: लंबाई के लिए
\emph{मीटर}\index{मीटर} (\unit{m}), द्रव्यमान के लिए
\emph{किलोग्राम}\index{किलोग्राम} (\unit{kg}) और समय के लिए
\emph{सेकंड}\index{सेकंड} (\unit{s})। (शेष चार --- ऐम्पियर, केल्विन, मोल,
कैंडेला --- तब जुड़ेंगे जब कहानी में विद्युत, ऊष्मा और रसायन आएँगे।)
\end{definition}

\begin{remark}[मात्रक के बिना संख्या कोई परिणाम नहीं]\label{rem:g10:orders-of-magnitude:units}
``दूरी $384$ है'' का कोई अर्थ नहीं: \qty{384}{mm}, \qty{384}{km} और \qty{384}{Mm} तीन अलग
दुनियाएँ हैं (आखिरी वाली संयोग से चंद्रमा तक की दूरी है)। 1999 में नासा
ने \$327 मिलियन का \emph{मार्स क्लाइमेट ऑर्बिटर} इसलिए खो दिया कि एक
प्रोग्राम ने प्रणोदक बलों की गणना पाउंड-बल में की और दूसरे ने उन्हें
न्यूटन में पढ़ा। इस पंक्ति से पुस्तक के अंत तक का घर का नियम: हर मापी
या गणना की गई संख्या अपना मात्रक साथ लिए चलती है, गणना की हर
पंक्ति पर।
\end{remark}

\begin{definition}[SI उपसर्ग]\label{def:g10:orders-of-magnitude:prefixes}
किसी मात्रक से जुड़ा \emph{उपसर्ग}\index{उपसर्ग (SI)} उसे दस की
एक निश्चित घात से गुणा कर देता है:
\begin{center}
\begin{tabular}{llc@{\qquad\qquad}llc}
फेम्टो & f & $10^{-15}$ & सेंटी & c & $10^{-2}$ \\
पीको  & p & $10^{-12}$ & किलो  & k & $10^{3}$  \\
नैनो  & n & $10^{-9}$  & मेगा  & M & $10^{6}$  \\
माइक्रो & $\mu$ & $10^{-6}$ & गीगा & G & $10^{9}$ \\
मिली & m & $10^{-3}$  & टेरा  & T & $10^{12}$ \\
\end{tabular}
\end{center}
इस प्रकार $\qty{1}{nm} = \qty{e-9}{m}$ और $\qty{1}{km} = \qty{e3}{m}$।
\end{definition}

\begin{example}[उपसर्ग पढ़ना]\label{ex:g10:orders-of-magnitude:prefixes}
हरे प्रकाश की तरंगदैर्घ्य लगभग $\qty{500}{nm} = \qty{5.00e-7}{m}$ होती है; मनुष्य का बाल लगभग
$\qty{70}{\micro m} = \qty{7.0e-5}{m}$ मोटा होता है; चंद्रमा $\qty{384}{Mm} = \qty{3.84e8}{m}$ पर परिक्रमा करता है। वर्षा की एक
बूँद का द्रव्यमान लगभग $\qty{30}{mg} = \qty{3.0e-5}{kg}$ होता है --- ध्यान दीजिए कि द्रव्यमान का
SI मूल मात्रक, सातों में अकेला, पहले से ही एक \omterm{def:g10:orders-of-magnitude:prefixes}{उपसर्ग}
लिए हुए है।
\end{example}

\section{वैज्ञानिक संकेतन और सार्थक अंक}

\begin{definition}[वैज्ञानिक संकेतन]\label{def:g10:orders-of-magnitude:scinot}
कोई संख्या \emph{वैज्ञानिक संकेतन}\index{वैज्ञानिक संकेतन} में तब होती
है जब उसे इस रूप में लिखा जाए:
\[
a \times 10^{n}, \qquad 1 \leq a < 10, \quad n \text{ एक पूर्णांक}
\]
गुणक $a$ \emph{अपूर्णांश}\index{अपूर्णांश} कहलाता है और पूर्णांक
$n$ \emph{घातांक}\index{घातांक}। इस प्रकार $\num{149600000} = \num{1.496e8}$ और $\num{0.00052} = \num{5.2e-4}$।
\end{definition}

\begin{proposition}[दस की घातों के साथ गणना]\label{prop:g10:orders-of-magnitude:powers}
सभी पूर्णांकों $m$ और $n$ के लिए:
\[
10^{m} \times 10^{n} = 10^{m+n}, \qquad
\frac{1}{10^{n}} = 10^{-n}, \qquad
\left(10^{m}\right)^{n} = 10^{mn}.
\]
\end{proposition}

\begin{proof}[आंशिक उपपत्ति]
धनात्मक घातांकों के लिए गुणनखंड गिन लीजिए: $10^{m} \times 10^{n}$ कुल $m + n$
दहाइयों का गुणनफल है, और $(10^m)^n$ में $n$ समूह, प्रत्येक $m$
दहाइयों का, एक के बाद एक जुड़ते हैं। $1/10^{n}$ के लिए $10^{-n}$ लिखना ठीक वही
परिपाटी है जो पहले नियम को सभी पूर्णांकों के लिए सच बनाए रखती है: वह
$10^{n} \times 10^{-n} = 10^{0} = 1$ को अनिवार्य कर देती है। साथी गणित का खंड \omterm{def:g10:orders-of-magnitude:scinot}{घातांक} के नियम
पूरे विस्तार से पढ़ता है।
\end{proof}

\begin{example}[वैज्ञानिक संकेतन में गुणा]\label{ex:g10:orders-of-magnitude:scinot}
द्रव्यमान के हिसाब से सूर्य में कितनी पृथ्वियाँ समाएँगी? $m_{\text{सूर्य}} = \qty{1.99e30}{kg}$ और
$m_{\text{पृथ्वी}} = \qty{5.97e24}{kg}$ लेकर:
\[
\frac{\num{1.99e30}}{\num{5.97e24}}
= \frac{1.99}{5.97} \times 10^{30-24}
= 0.333 \times 10^{6}
= \num{3.33e5}.
\]
अपूर्णांशों और घातांकों को अलग-अलग सँभाला जाता है; अंतिम
उत्तर को फिर से सामान्य कर दिया जाता है ताकि \omterm{def:g10:orders-of-magnitude:scinot}{अपूर्णांश} वापस
$[1, 10)$ में आ जाए।
\end{example}

\begin{definition}[सार्थक अंक]\label{def:g10:orders-of-magnitude:sigfig}
किसी मापी गई संख्या के \emph{सार्थक अंक}\index{सार्थक अंक} वे अंक हैं
जिनकी माप वास्तव में गवाही देती है: आरंभ के शून्यों को छोड़कर सभी लिखे
गए अंक। अंत के शून्य गिने जाते हैं --- $\qty{5.20e-3}{s}$ (तीन सार्थक अंक) $\qty{5.2e-3}{s}$
(दो) से अधिक का वादा करता है। \omterm{def:g10:orders-of-magnitude:scinot}{वैज्ञानिक संकेतन} यह गिनती एक ही
नज़र में दिखा देता है: \omterm{def:g10:orders-of-magnitude:scinot}{अपूर्णांश} ठीक उतने ही अंक लिए होता है
जितने सार्थक हैं।
\end{definition}

\begin{method}[कितने अंक रखें]\label{met:g10:orders-of-magnitude:sigfig}
कोई परिकलित परिणाम अपनी सामग्री से अधिक परिशुद्ध नहीं हो सकता।
\begin{enumerate}
  \item गुणनफल या भागफल में उतने ही \omterm{def:g10:orders-of-magnitude:sigfig}{सार्थक अंक} रखिए जितने
        \emph{सबसे कम} परिशुद्ध गुणक में हैं।
  \item केवल अंतिम उत्तर को पूर्णांकित कीजिए, बीच के चरणों को कभी नहीं।
  \item संदेह हो तो तीन \omterm{def:g10:orders-of-magnitude:sigfig}{सार्थक अंक} --- इस पुस्तक की कामकाजी
        परिशुद्धि यही है।
\end{enumerate}
\end{method}

\begin{example}[पृथ्वी की परिधि]\label{ex:g10:orders-of-magnitude:sigfig}
पृथ्वी का व्यास $d = \qty{1.27e7}{m}$ है (तीन \omterm{def:g10:orders-of-magnitude:sigfig}{सार्थक अंक}), इसलिए उसकी परिधि
$\pi d = \num{3.14159} \times \qty{1.27e7}{m} = \qty{3.99e7}{m}$ है --- तीन अंक, न कि आठ, जो कैलकुलेटर दिखाता है। $\qty{39898227}{m}$ लिखने का
अर्थ होता ग्रह को \omterm{def:g10:orders-of-magnitude:unit}{मीटर} तक जानने का दावा करना।
\end{example}

\begin{omfigure}
\begin{tikzpicture}[x=1.9cm, font=\small]
  \draw[-Stealth] (-0.3,0) -- (5.45,0);
  \foreach \e in {0,...,5} {
    \draw (\e,-0.1) -- (\e,0.1);
    \node[below] at (\e,-0.12) {$10^{\e}$};
  }
  \foreach \x in {2.301,2.477,2.602,2.699,2.778,2.845,2.903,2.954}
    \draw (\x,-0.05) -- (\x,0.05);
  \draw[Stealth-Stealth, omMeth] (0,0.45) -- (1,0.45)
    node[midway, above] {$\times 10$};
  \fill[omDef] (2.477,0) circle (1.7pt);
  \node[above=3pt, omDef] at (2.477,0.06) {$3 \times 10^{2}$};
  \draw[-Stealth, omDef, dashed] (2.38,0.22) -- (2.06,0.22);
  \fill[omThm] (3.845,0) circle (1.7pt);
  \node[above=3pt, omThm] at (3.845,0.06) {$7 \times 10^{3}$};
  \draw[-Stealth, omThm, dashed] (3.90,0.22) -- (3.96,0.22);
\end{tikzpicture}

{\small एक ऐसी मापनी जिस पर हर बराबर कदम दस से \emph{गुणा} करता है।
छोटी लकीरें $200, 300, \dots, 900$ हैं: पूरा एक दशक हर खंड में ठुँस जाता है। इस मापनी
पर $\num{300}$ $10^{2}$ की ओर झुकता है और $\num{7000}$ $10^{4}$ की ओर --- यही उनकी
परिमाण की कोटियाँ हैं।}
\end{omfigure}

\section{ब्रह्मांड की सीढ़ी}

\begin{definition}[परिमाण की कोटि]\label{def:g10:orders-of-magnitude:oom}
किसी राशि की \emph{परिमाण की कोटि}\index{परिमाण की कोटि} दस की वह घात
है जो उसके सबसे पास पड़ती है: राशि को $a \times 10^{n}$ के रूप में लिखिए, जहाँ
$1 \leq a < 10$, और $a < 5$ हो तो $10^{n}$ लीजिए, $a \geq 5$ हो तो $10^{n+1}$। (सीमा की
ठीक-ठीक परिपाटी कभी मायने नहीं रखती: परिमाण की कोटि का कथन दस के गुणक
तक ही सही माना जाता है।)
\end{definition}

\begin{example}[परिमाण की कोटियाँ]\label{ex:g10:orders-of-magnitude:oom}
मनुष्य का जीवन लगभग $80$ वर्ष चलता है, यानी $80 \times \qty{3.16e7}{s} \approx \qty{2.5e9}{s}$: \omterm{def:g10:orders-of-magnitude:oom}{परिमाण की
कोटि} $\qty{e9}{s}$ --- कुछ अरब \omterm{def:g10:orders-of-magnitude:unit}{सेकंड}, इन्हें ठीक से बिताइए। मनुष्य $\qty{1.7}{m}$
लंबा होता है: कोटि $\qty{e0}{m}$। माउंट एवरेस्ट, $\qty{8848}{m}$: कोटि $\qty{e4}{m}$,
क्योंकि $8.848 \geq 5$।
\end{example}

ऐसे अनुमानों को \omterm{def:g10:orders-of-magnitude:scinot}{घातांक} के हिसाब से एक के ऊपर एक रखने से इस
अध्याय का मुख्य चित्र बनता है: एक सीढ़ी, जिसका हर डंडा नीचे वाले से दस
गुना चौड़ी दुनिया है।

\begin{omfigure}
\begin{tikzpicture}[x=0.26cm, font=\small]
  \draw[-Stealth] (-16.5,0) -- (28.7,0);
  \foreach \e in {-15,-10,-5,0,5,10,15,20,25} {
    \draw (\e,-0.12) -- (\e,0.12);
    \node[below, font=\footnotesize] at (\e,-0.14) {$10^{\e}$};
  }
  \foreach \e/\t in {-15/नाभिक, -10/परमाणु, -6/जीवाणु, 0/मनुष्य,
                     7/पृथ्वी, 11/पृथ्वी--सूर्य, 16/प्रकाश वर्ष,
                     21/आकाशगंगा, 27/ब्रह्मांड} {
    \draw[omDef, thick] (\e,0.14) -- (\e,0.62);
    \node[rotate=52, anchor=south west, inner sep=1pt, omDef]
      at (\e,0.66) {\t};
  }
  \node[below, font=\footnotesize] at (26.6,-0.75) {आकार (m)};
\end{tikzpicture}

{\small ब्रह्मांड की सीढ़ी, मीटरों में: परमाणु नाभिक और
प्रेक्षणीय ब्रह्मांड के बीच दस की $42$ घातें हैं, और मनुष्य परमाणु
से दस डंडे ऊपर खड़ा है, पर चोटी से सत्ताईस नीचे।}
\end{omfigure}

पृथ्वी वाले डंडे के ऊपर \omterm{def:g10:orders-of-magnitude:unit}{मीटर} सुविधाजनक नहीं रह जाते; खगोल
विज्ञान अपने मात्रक रखता है, और दोनों इसी सीढ़ी से परिभाषित
होते हैं।

\begin{definition}[खगोलीय मात्रक और प्रकाश वर्ष]\label{def:g10:orders-of-magnitude:lightyear}
\emph{खगोलीय मात्रक}\index{खगोलीय मात्रक} ($\qty{1}{au} = \qty{1.496e11}{m}$) पृथ्वी और सूर्य के
बीच की औसत दूरी है। \emph{प्रकाश वर्ष}\index{प्रकाश वर्ष} (\unit{ly}) वह
दूरी है जो प्रकाश एक वर्ष में तय करता है। निर्वात में प्रकाश
\[
c = \qty{3.00e8}{m/s},
\]
की चाल से चलता है, जो प्रकृति की अनुमत सबसे तेज़ चाल है --- इसलिए
प्रकाश वर्ष एक दूरी है, समय नहीं।
\end{definition}

\begin{omfigure}
\begin{tikzpicture}[x=1.08cm, font=\small]
  \draw[-Stealth] (7.6,0) -- (17.5,0);
  \foreach \e in {8,10,12,14,16} {
    \draw (\e,-0.12) -- (\e,0.12);
    \node[below, font=\footnotesize] at (\e,-0.14) {$10^{\e}$};
  }
  \foreach \x/\t/\d/\a in {8.58/चंद्रमा/{\qty{1.3}{s}}/north,
                           11.17/सूर्य/{\qty{8.3}{min}}/north,
                           12.65/नेपच्यून/{\qty{4.2}{h}}/north,
                           15.98/{$\qty{1}{ly}$}/{\qty{1}{yr}}/north east,
                           16.60/प्रॉक्सिमा/{\qty{4.2}{yr}}/north west} {
    \draw[omDef, thick] (\x,0.14) -- (\x,0.55);
    \node[rotate=52, anchor=south west, inner sep=1pt, omDef]
      at (\x,0.58) {\t};
    \node[anchor=\a, font=\footnotesize, omThm] at (\x,-0.55) {\d};
  }
  \node[right, font=\footnotesize] at (17.5,0) {दूरी (m)};
\end{tikzpicture}

{\small गुणात्मक मापनी पर पृथ्वी से दूरियाँ, और हर एक के नीचे लाल रंग
में प्रकाश का यात्रा-समय: चंद्रमा प्रकाश-सेकंड दूर है, ग्रह
प्रकाश-घंटे, और तारे \omterm{def:g10:orders-of-magnitude:lightyear}{प्रकाश वर्ष}।}
\end{omfigure}

\begin{example}[प्रकाश वर्ष, मीटरों में]\label{ex:g10:orders-of-magnitude:ly}
एक वर्ष में $365.25 \times 24 \times 3600 = \qty{3.156e7}{s}$ होते हैं (जो सुखद रूप से $\pi \times 10^{7}\,\unit{s}$ के पास है)। अतः
\[
\qty{1}{ly} = c \times \qty{3.156e7}{s}
= \qty{3.00e8}{m/s} \times \qty{3.156e7}{s}
= \qty{9.47e15}{m}.
\]
सूर्य के बाद का सबसे पास का तारा प्रॉक्सिमा सेंटौरी $\qty{4.2}{ly} \approx \qty{4.0e16}{m}$ दूर है:
प्रकाश, जो पृथ्वी--सूर्य की दूरी आठ मिनट में पार कर लेता है, वहाँ से आने
में चार साल से अधिक लेता है।
\end{example}

\begin{method}[परिमाण की कोटि से जाँच]\label{met:g10:orders-of-magnitude:sanity}
किसी भी परिकलित परिणाम पर भरोसा करने से पहले:
\begin{enumerate}
  \item हर आगत को उसकी \omterm{def:g10:orders-of-magnitude:oom}{परिमाण की कोटि} तक पूर्णांकित कीजिए;
  \item \cref{prop:g10:orders-of-magnitude:powers} की सहायता से गणना केवल दस की घातों से दोहराइए ---
        इसमें कुछ ही \omterm{def:g10:orders-of-magnitude:unit}{सेकंड} लगते हैं;
  \item तुलना कीजिए: परिशुद्ध उत्तर को अनुमान से दस के गुणक भर के भीतर
        मेल खाना चाहिए, और उसका मात्रक सही होना चाहिए। इनमें
        से कोई भी जाँच बिगड़े तो आगे बढ़ने से पहले गलती खोज निकालिए।
\end{enumerate}
\end{method}

\begin{remark}[इसकी ज़रूरत क्यों]\label{rem:g10:orders-of-magnitude:sanity}
कैलकुलेटर खुशी-खुशी बता देगा कि किसी सेब का द्रव्यमान $\qty{4e8}{kg}$ है ---
उसे इससे बेहतर पता ही नहीं हो सकता। कोई सूत्र बेतुकेपन पर झंडी नहीं
उठाता; आदत उठाती है। भौतिक विज्ञानी पहले अनुमान लगाते हैं, गणना बाद में
करते हैं: अनुमान खिसके हुए दशमलव बिंदु, अदल-बदल गए मात्रक और
उलटी हो गई भिन्नें पकड़ लेता है --- हर स्तर पर अधिकांश गलत उत्तरों के
पीछे यही तीन गलतियाँ होती हैं।
\end{remark}

\section{मापन और अनिश्चितता}

कोई उपकरण सही मान नहीं पढ़ता। कोई माप तभी ईमानदार है जब वह यह भी बताए
कि वह कितनी दूर तक चूक सकता है।

\begin{definition}[निरपेक्ष और सापेक्ष अनिश्चितता]\label{def:g10:orders-of-magnitude:uncertainty}
किसी राशि $x$ की माप इस तरह बताई जाती है:
\[
x = x_{\text{m}} \pm \Delta x ,
\]
जहाँ $x_{\text{m}}$ मापा गया मान है और उसी मात्रक में दी गई
\emph{निरपेक्ष अनिश्चितता}\index{अनिश्चितता!निरपेक्ष} $\Delta x > 0$ संभावित
त्रुटि को सीमित करती है: दावा यह है कि सही मान $x_{\text{m}} - \Delta x$ और $x_{\text{m}} + \Delta x$ के
बीच है। \emph{सापेक्ष अनिश्चितता}\index{अनिश्चितता!सापेक्ष} अनुपात
$\Delta x / x_{\text{m}}$ है, जिसे प्रायः प्रतिशत में बताया जाता है; ``परिशुद्ध'' शब्द
इसी को नापता है।
\end{definition}

\begin{method}[अंशांकित उपकरण पढ़ना]\label{met:g10:orders-of-magnitude:reading}
मापनी, मापक सिलिंडर या डायल के लिए:
\begin{enumerate}
  \item निशान के सबसे पास वाला अंशांक पढ़िए, आँख पैमाने के लंबवत रखकर
        (तिरछी आँख पढ़त को खिसका देती है);
  \item सबसे छोटे अंशांक का कम से कम आधा \omterm{def:g10:orders-of-magnitude:uncertainty}{निरपेक्ष अनिश्चितता}
        मानिए --- और अधिक, यदि वस्तु का किनारा या द्रव की सतह धुँधली
        हो;
  \item परिणाम को उसके मात्रक के साथ $x_{\text{m}} \pm \Delta x$ के रूप में
        लिखिए, और अनिश्चित अंक से आगे कोई अंक मत रखिए।
\end{enumerate}
\end{method}

\begin{omfigure}
\begin{tikzpicture}[x=0.9cm, font=\small]
  \fill[omThm, opacity=0.15] (6.5,-0.85) rectangle (7.5,1.45);
  \draw (-1.2,0) rectangle (10.6,1.1);
  \foreach \x in {0,...,10} \draw (\x,1.1) -- (\x,0.72);
  \foreach \x/\l in {0/120, 5/125, 10/130}
    \node[font=\footnotesize] at (\x,0.4) {\l};
  \fill[gray!25] (-1.0,-0.75) rectangle (7.0,-0.12);
  \draw (-1.0,-0.75) rectangle (7.0,-0.12);
  \draw[dashed, omThm] (7,-0.85) -- (7,1.45);
  \node[anchor=west, omThm] at (7.7,-0.5)
    {$L = (127.0 \pm 0.5)\,mm$};
\end{tikzpicture}

{\small मिलीमीटर मापनी पढ़ना: छड़ का सिरा $\qty{127}{mm}$ के अंशांक के सबसे पास
पड़ता है, और दोनों ओर आधे-आधे अंशांक की छायांकित पट्टी ही ईमानदार दावा
है।}
\end{omfigure}

\begin{example}[सापेक्ष परिशुद्धि]\label{ex:g10:orders-of-magnitude:reading}
ऊपर की पढ़त की \omterm{def:g10:orders-of-magnitude:uncertainty}{सापेक्ष अनिश्चितता} $\frac{0.5}{127} \approx 0.4\%$ है। वही मापनी जब
$\qty{12}{mm}$ का बोल्ट मापती है तो $\frac{0.5}{12} \approx 4\%$ देती है: वही उपकरण छोटी वस्तु पर
दस गुना कम परिशुद्ध है। परिशुद्धि \emph{माप} का गुण है, औज़ार का नहीं।
\end{example}

\begin{remark}[छोटी निरपेक्ष, विशाल सापेक्ष --- और इसका उलटा]\label{rem:g10:orders-of-magnitude:precision}
लेज़र केंद्र पृथ्वी--चंद्रमा की दूरी $\qty{3.84e8}{m}$ को लगभग $\pm\qty{3}{cm}$ तक मापते
हैं: $10^{-10}$ की \omterm{def:g10:orders-of-magnitude:uncertainty}{सापेक्ष अनिश्चितता}, जो अब तक की सबसे पैनी
मापों में गिनी जाती है --- और वह भी ऐसी निरपेक्ष त्रुटि के साथ जिस पर
कोई मापनी इतराएगी नहीं। हमेशा पूछिए, \emph{किसके सापेक्ष}: पूरी कहानी
युग्म $(\Delta x,\ \Delta x / x_{\text{m}})$ बताता है, और इनमें से कोई एक संख्या अकेले भटका सकती है।
\end{remark}

\section{अभ्यास}

\begin{exercise}[$\star$]\label{exo:g10:orders-of-magnitude:1}
\omterm{def:g10:orders-of-magnitude:scinot}{वैज्ञानिक संकेतन} में, मीटरों में बदलिए:
\begin{enumerate}
  \item $\qty{750}{nm}$ (लाल प्रकाश की तरंगदैर्घ्य);
  \item $\qty{2.5}{km}$;
  \item $\qty{3.2}{\micro m}$ (मिट्टी का एक जीवाणु);
  \item $\qty{0.15}{Tm}$ (पृथ्वी--सूर्य की दूरी, एक असामान्य पोशाक में)।
\end{enumerate}
\end{exercise}

\begin{exercise}[$\star$]\label{exo:g10:orders-of-magnitude:2}
\omterm{def:g10:orders-of-magnitude:scinot}{वैज्ञानिक संकेतन} में लिखिए और सार्थक अंकों की संख्या
बताइए:
\begin{enumerate}
  \item पृथ्वी--सूर्य की दूरी, $\qty{149600000}{km}$, मीटरों में;
  \item हाइड्रोजन परमाणु का व्यास, $\qty{0.000000000106}{m}$;
  \item विश्व की जनसंख्या, लगभग $\num{8100000000}$;
  \item $\qty{0.00520}{s}$ की एक अवधि।
\end{enumerate}
\end{exercise}

\begin{exercise}[$\star$]\label{exo:g10:orders-of-magnitude:3}
कैलकुलेटर के बिना गणना कीजिए और हर परिणाम \omterm{def:g10:orders-of-magnitude:scinot}{वैज्ञानिक संकेतन} में
दीजिए:
\[
(3\times10^{8})\times(2\times10^{-5}), \qquad
\frac{8\times10^{4}}{4\times10^{-3}}, \qquad
(5\times10^{6})^{2}, \qquad
\frac{(4\times10^{-3})\times(3\times10^{-9})}{6\times10^{-7}} .
\]
\end{exercise}

\begin{exercise}[$\star$]\label{exo:g10:orders-of-magnitude:4}
इनकी \omterm{def:g10:orders-of-magnitude:oom}{परिमाण की कोटि} उसके मात्रक सहित बताइए: $\qty{1.7}{m}$
की मानव-लंबाई; माउंट एवरेस्ट, $\qty{8848}{m}$; एक \emph{ई.~कोलाई} जीवाणु,
$\qty{2e-6}{m}$; पृथ्वी का द्रव्यमान, $\qty{5.97e24}{kg}$।
\end{exercise}

\begin{exercise}[$\star$]\label{exo:g10:orders-of-magnitude:5}
एक मापनी मिलीमीटर में अंशांकित है। आप एक पेंसिल को $\qty{127}{mm}$ मापते हैं।
\begin{enumerate}
  \item परिणाम को उसकी \omterm{def:g10:orders-of-magnitude:uncertainty}{निरपेक्ष अनिश्चितता} सहित लिखिए।
  \item \omterm{def:g10:orders-of-magnitude:uncertainty}{सापेक्ष अनिश्चितता} निकालिए।
  \item वही मापनी $\qty{12}{mm}$ की रबड़ मापती है। दोनों मापों में कौन-सी अधिक
        परिशुद्ध है, और कितने गुना?
\end{enumerate}
\end{exercise}

\begin{exercise}[$\star\star$]\label{exo:g10:orders-of-magnitude:6}
$c = \qty{3.00e8}{m/s}$ और $\qty{1}{ly} = \qty{9.47e15}{m}$ लेकर:
\begin{enumerate}
  \item प्रॉक्सिमा सेंटौरी ($\qty{4.24}{ly}$) और सिरियस ($\qty{8.6}{ly}$) की दूरियाँ
        मीटरों में व्यक्त कीजिए।
  \item यान \emph{वॉयेजर 1} पृथ्वी से लगभग $\qty{2.4e13}{m}$ दूर है। $c$ की
        चाल से चलते हुए उसके रेडियो संकेत को हम तक पहुँचने में कितना
        समय लगता है? उत्तर घंटों में दीजिए।
\end{enumerate}
\end{exercise}

\begin{exercise}[$\star\star$]\label{exo:g10:orders-of-magnitude:7}
चंद्रमा ($\qty{3.84e8}{m}$), सूर्य ($\qty{1.496e11}{m}$) और नेपच्यून ($\qty{4.50e12}{m}$) से पृथ्वी तक
प्रकाश को पहुँचने में कितना समय लगता है? हर उत्तर समझदारी से चुने गए
मात्रक में दीजिए, और वाक्य पूरा कीजिए: ``जब आप नेपच्यून को
देखते हैं, तो आप उसे वैसा देखते हैं जैसा वह\dots''
\end{exercise}

\begin{exercise}[$\star\star$]\label{exo:g10:orders-of-magnitude:8}
कागज़ की $500$ शीटों की एक गड्डी मिलीमीटर मापनी से मापने पर $\qty{52.0}{mm}$
मोटी निकलती है, अनिश्चितता $\pm\qty{0.5}{mm}$ के साथ।
\begin{enumerate}
  \item एक शीट की मोटाई उसकी \omterm{def:g10:orders-of-magnitude:uncertainty}{निरपेक्ष अनिश्चितता} सहित निकालिए।
  \item उस मोटाई की \omterm{def:g10:orders-of-magnitude:uncertainty}{सापेक्ष अनिश्चितता} क्या है?
  \item उसी मापनी से एक अकेली शीट मापने के बजाय पूरी गड्डी मापना बहुत
        बेहतर क्यों है? संख्याओं में बताइए।
\end{enumerate}
\end{exercise}

\begin{exercise}[$\star\star$]\label{exo:g10:orders-of-magnitude:9}
परमाणु का व्यास $\qty{2e-10}{m}$ लीजिए। अगल-बगल रखे कितने परमाणु इन्हें पाट
देंगे: $\qty{2.0}{\micro m}$ का एक जीवाणु; $\qty{70}{\micro m}$ का एक बाल; $\qty{1.7}{m}$ का एक मनुष्य?
आखिरी उत्तर \omterm{def:g10:orders-of-magnitude:oom}{परिमाण की कोटि} के रूप में भी दीजिए।
\end{exercise}

\begin{exercise}[$\star\star$]\label{exo:g10:orders-of-magnitude:10}
एक मापनी-प्रतिरूप में सूर्य (व्यास $\qty{1.39e9}{m}$) को $\qty{22}{cm}$ व्यास वाली एक
फुटबॉल से दिखाया गया है।
\begin{enumerate}
  \item प्रतिरूप का मापनी-गुणक निकालिए।
  \item पृथ्वी ($\qty{1.27e7}{m}$) का प्रतिरूप-व्यास और फुटबॉल से उसकी
        प्रतिरूप-दूरी ज्ञात कीजिए (वास्तविक दूरी $\qty{1.496e11}{m}$)।
  \item प्रतिरूप-पृथ्वी से प्रतिरूप-चंद्रमा कितनी दूर है (वास्तविक दूरी
        $\qty{3.84e8}{m}$)? और प्रतिरूप-नेपच्यून कितनी दूर है (वास्तविक दूरी
        $\qty{4.50e12}{m}$)?
\end{enumerate}
\end{exercise}

\begin{exercise}[$\star\star$]\label{exo:g10:orders-of-magnitude:11}
\omterm{def:g10:orders-of-magnitude:sigfig}{सार्थक अंक} काम पर।
\begin{enumerate}
  \item एक A4 शीट $\qty{21.0}{cm}$ गुणा $\qty{29.7}{cm}$ की है। उसका क्षेत्रफल
        सार्थक अंकों की सही संख्या के साथ दीजिए।
  \item एक धावक $\qty{100.0}{m}$ की दूरी $\qty{12.3}{s}$ में तय करता है। औसत चाल सही
        पूर्णांकन के साथ दीजिए।
  \item एक वाक्य में बताइए कि क्षेत्रफल को $\qty{623.70}{cm^2}$ बताने में क्या
        बेईमानी होगी।
\end{enumerate}
\end{exercise}

\begin{exercise}[$\star\star\star$]\label{exo:g10:orders-of-magnitude:12}
$\qty{1.0}{mm^3}$ आयतन के तेल की एक बूँद शांत पानी पर डाली जाती है और $\qty{1.1}{m}$
व्यास की एक गोल परत में फैल जाती है --- ऐसी परत जो, आगे के अध्याय
दिखाएँगे, एक अणु जितनी मोटी है।
\begin{enumerate}
  \item परत का क्षेत्रफल निकालिए, फिर उसकी मोटाई।
  \item इससे एक अणु के आकार की \omterm{def:g10:orders-of-magnitude:oom}{परिमाण की कोटि} निकालिए, और उसे
        सीढ़ी के परमाणु वाले डंडे से मिलाकर जाँचिए।
\end{enumerate}
(यह मोटा-मोटी अनुमान वाला प्रयोग सबसे पहले बेंजामिन फ्रैंकलिन ने लंदन
के एक तालाब पर किया था, और यह अणुओं की दुनिया की मानवता की पहली ईमानदार
मापों में से एक था।)
\end{exercise}

\begin{exercise}[$\star\star\star$]\label{exo:g10:orders-of-magnitude:13}
मनुष्य के पूरे जीवन में धड़कनों की संख्या का अनुमान लगाइए। अपनी
मान्यताएँ बताइए (विश्राम की हृदय-गति, आयु), अंकगणित को दस की घातों और
\omterm{def:g10:orders-of-magnitude:scinot}{अपूर्णांश} के एक ही अंक तक सीमित रखिए, और अंतिम उत्तर केवल
\omterm{def:g10:orders-of-magnitude:oom}{परिमाण की कोटि} के रूप में दीजिए। इसे चार सार्थक अंकों
तक बताना निरर्थक क्यों होगा?
\end{exercise}

\begin{exercise}[$\star\star\star$]\label{exo:g10:orders-of-magnitude:14}
कागज़ की एक शीट $\qty{0.10}{mm}$ मोटी है, और हर तह ढेर की मोटाई दुगुनी कर देती
है। $2^{10} \approx 10^{3}$ लेकर:
\begin{enumerate}
  \item कितनी तहों के बाद ढेर माउंट एवरेस्ट ($\qty{8848}{m}$) से ऊपर निकल
        जाएगा?
  \item कितनी तहों के बाद वह चंद्रमा ($\qty{3.84e8}{m}$) तक पहुँच जाएगा?
  \item व्यवहार में एक शीट को लगभग सात बार से अधिक नहीं मोड़ा जा सकता।
        क्या इससे अंकगणित गलत हो जाता है, या केवल प्रयोग?
\end{enumerate}
\end{exercise}

\begin{exercise}[$\star\star\star$]\label{exo:g10:orders-of-magnitude:15}
लेज़र परासन से पृथ्वी--चंद्रमा की दूरी $\qty{3.84e8}{m} \pm \qty{3}{cm}$ निकलती है; एक अच्छी
मापनी किसी मेज़ को $\qty{1.000}{m} \pm \qty{1}{mm}$ बताती है।
\begin{enumerate}
  \item दोनों सापेक्ष अनिश्चितताएँ निकालिए। कौन-सी माप अधिक परिशुद्ध
        है, और कितने गुना?
  \item चंद्रमा वाली माप की सापेक्ष परिशुद्धि से मेल खाने के लिए आपको
        मेज़ को किस \omterm{def:g10:orders-of-magnitude:uncertainty}{निरपेक्ष अनिश्चितता} तक मापना पड़ेगा? उस
        लंबाई की तुलना परमाणु के आकार से कीजिए, और उचित आदर के साथ
        निष्कर्ष निकालिए।
\end{enumerate}
\end{exercise}

\section{समस्या: परमाणु से एंड्रोमेडा तक}

\begin{problem}[{सप्ताहांत समस्या --- रेत का एक अकेला कण परमाणु, सौर
मंडल, आकाशगंगा और प्रेक्षणीय ब्रह्मांड को नाप देता है, और इनके बीच की
सीढ़ी पर आपकी जगह भी बता देता है}]
\label{pb:g10:orders-of-magnitude:1}
रेत का कण भौतिकी की सबसे मामूली वस्तु है, और इस सप्ताहांत वह अतिरिक्त
काम करेगा: पहले हम उसके परमाणु गिनेंगे, फिर सूर्य को उसी आकार तक सिकोड़
कर ब्रह्मांड को कदमों में नापेंगे, फिर सर्वेक्षण का काम प्रकाश को
सौंपेंगे, और अंत में \emph{आप} में परमाणु गिनेंगे। पूरे समय काम आने
वाले आँकड़े: कण का व्यास $\qty{0.5}{mm}$; परमाणु का व्यास $\qty{2.5e-10}{m}$; रेत का घनत्व
$\qty{2500}{kg/m^3}$; व्यास: सूर्य $\qty{1.39e9}{m}$, पृथ्वी $\qty{1.27e7}{m}$, आकाशगंगा $\qty{9.5e20}{m}$,
प्रेक्षणीय ब्रह्मांड $\qty{8.8e26}{m}$; हमसे दूरियाँ: चंद्रमा $\qty{3.84e8}{m}$, सूर्य
$\qty{1.496e11}{m}$, नेपच्यून $\qty{4.5e12}{m}$, प्रॉक्सिमा सेंटौरी $\qty{4.0e16}{m}$, एंड्रोमेडा
आकाशगंगा $\qty{2.4e22}{m}$; $c = \qty{3.00e8}{m/s}$; पानी के एक अणु का द्रव्यमान $\qty{3.0e-26}{kg}$ है और
उसमें तीन परमाणु होते हैं।

\medskip
\noindent\textbf{भाग I --- कण की बनावट।}
\begin{enumerate}
  \item कण का व्यास और परमाणु का व्यास \omterm{def:g10:orders-of-magnitude:scinot}{वैज्ञानिक संकेतन} में,
        मीटरों में लिखिए, और बताइए कि हर एक किस SI
        \omterm{def:g10:orders-of-magnitude:prefixes}{उपसर्ग} के सबसे पास है।
  \item एक कण के व्यास भर में अगल-बगल कितने परमाणु बैठते हैं?
  \item कण को $\qty{0.5}{mm}$ भुजा के ऐसे घन के रूप में लीजिए जो परमाणुओं से
        ठसाठस भरा है। कण में परमाणुओं की संख्या का अनुमान
        \omterm{def:g10:orders-of-magnitude:oom}{परिमाण की कोटि} के रूप में लगाइए।
  \item उस घन का आयतन निकालिए, फिर कण का द्रव्यमान। द्रव्यमान को
        मिलीग्राम में बताइए।
  \item प्रतिरूप की जाँच कीजिए: कण के द्रव्यमान को उसके परमाणुओं की
        संख्या से भाग दीजिए और परमाणु-द्रव्यमानों की ज्ञात
        \omterm{def:g10:orders-of-magnitude:oom}{परिमाण की कोटि} (कुछ $10^{-26}\,\unit{kg}$) से तुलना कीजिए। फैसला?
\end{enumerate}

\medskip
\noindent\textbf{भाग II --- सूर्य, रेत के एक कण जितना।}
हम ब्रह्मांड का एक ऐसा मापनी-प्रतिरूप बनाते हैं जिसमें सूर्य वही कण है:
हर प्रतिरूप-आकार वास्तविक आकार को एक गुणक $s$ से गुणा करने पर
मिलता है।
\begin{enumerate}[resume]
  \item मापनी-गुणक $s$ निकालिए।
  \item प्रतिरूप में पृथ्वी का व्यास और कण-सूर्य से उसकी दूरी ज्ञात
        कीजिए। भाग I की दुनिया की कौन-सी वस्तु मोटे तौर पर
        प्रतिरूप-पृथ्वी जितनी है?
  \item प्रतिरूप-नेपच्यून कण से कितनी दूर परिक्रमा करता है? कौन-सा
        फर्नीचर पूरे ग्रह-मंडल को समा लेगा?
  \item प्रतिरूप-प्रॉक्सिमा सेंटौरी कितनी दूर है?
  \item रेत के दो कण, $\qty{14}{km}$ की दूरी पर, और बीच में लगभग कुछ भी नहीं:
        एक-दो वाक्यों में बताइए कि यह प्रतिरूप आकाशगंगा के खालीपन के
        बारे में क्या खोल देता है।
\end{enumerate}

\medskip
\noindent\textbf{भाग III --- सर्वेक्षण प्रकाश करता है।}
\begin{enumerate}[resume]
  \item रेत के असली कण को पार करने में प्रकाश को कितना समय लगता है?
  \item चंद्रमा से और सूर्य से हम तक पहुँचने में प्रकाश को कितना समय
        लगता है?
  \item $c$ और वर्ष की लंबाई से \omterm{def:g10:orders-of-magnitude:lightyear}{प्रकाश वर्ष} को
        मीटरों में फिर से निकालिए, और जाँचिए कि आँकड़ों वाली
        प्रॉक्सिमा दूरी लगभग $\qty{4.2}{ly}$ बैठती है।
  \item एंड्रोमेडा की दूरी को प्रकाश वर्षों में बदलिए। जो
        प्रकाश इस समय आपकी आँख में घुस रहा है, वह मोटे तौर पर कब
        एंड्रोमेडा से चला था, और उस समय पृथ्वी पर कौन चल रहा था?
  \item प्रकाश एक नैनोसेकंड में कितनी दूरी तय करता है? $\qty{3}{GHz}$ पर चलने
        वाला प्रोसेसर हर एक-तिहाई नैनोसेकंड में एक क्रिया शुरू करता है:
        एक वाक्य में बताइए कि तेज़ कंप्यूटर का छोटा होना क्यों ज़रूरी
        है।
\end{enumerate}

\medskip
\noindent\textbf{भाग IV --- आप, परमाणु और एंड्रोमेडा के बीच।}
\begin{enumerate}[resume]
  \item मनुष्य मोटे तौर पर $\qty{70}{kg}$ पानी है। आँकड़ों की सहायता से पानी
        में एक परमाणु के औसत द्रव्यमान का अनुमान लगाइए, फिर मनुष्य में
        परमाणुओं की संख्या का।
  \item उस गिनती की तुलना कण के परमाणुओं से (प्रश्न 3) और प्रेक्षणीय
        ब्रह्मांड के लगभग $2\times10^{23}$ तारों से कीजिए। दोनों गुणक दीजिए।
  \item सीढ़ी पर आप परमाणु के पास हैं या प्रेक्षणीय ब्रह्मांड के?
        अपने लिए $\qty{1.7}{m}$ लेकर अनुपात $\text{(आप)}/\text{(परमाणु)}$ और $\text{(ब्रह्मांड)}/\text{(आप)}$ की तुलना
        कीजिए।
  \item वह आकार $x$ ज्ञात कीजिए जो नाभिक ($\qty{e-15}{m}$) और प्रेक्षणीय
        ब्रह्मांड के बीच सीढ़ी पर ठीक आधे पर बैठता है: $x$ को
        $x / 10^{-15} = 10^{26.9} / x$ संतुष्ट करना चाहिए ($\qty{8.8e26}{m} \approx 10^{26.9}\,\unit{m}$ लेकर)। सौर मंडल की कौन-सी
        वस्तु संयोग से ठीक इसी आकार की है?
  \item और अंतिम चोट। प्रेक्षणीय ब्रह्मांड को भाग II वाले गुणक $s$
        से सिकोड़िए, ताकि सूर्य रेत का एक कण बन जाए। प्रतिरूप-ब्रह्मांड
        का व्यास मीटरों में दीजिए, फिर खगोलीय
        मात्रकों में, और एक वाक्य में निष्कर्ष निकालिए: हर तारे को रेत
        का कण बना देने पर भी ब्रह्मांड हर मानवीय पैमाने से छलक जाता है
        --- उसे केवल दस की घातें थाम पाती हैं।
\end{enumerate}
\end{problem}
